% migrate-to: futures-execution
\subsection{Futures Market Microstructure}

\begin{principle} \hlt{Centralised Futures Trading}\\
Liquid futures are highly electronic and usually trade on centralised limit order books. Fragmentation is limited because margining, clearing, delivery rules, and contract fungibility make post-trade duplication costly.
\end{principle}

\begin{caveat} \hlt{No Call Auction Norm}\\
Many futures markets do not rely on equity-style opening and closing call auctions. Instead, participants agree on liquid trading hours during which the contract can support reliable price discovery and hedging.
\end{caveat}

\begin{mechanism} \hlt{Lead-Lag With Cash Markets}\\
Index futures and their cash baskets form prices jointly.\\
Futures may lead when they are cheaper to trade, more leveraged, or open during information arrival; cash markets may lead when constituent-level information dominates.
\end{mechanism}

\begin{principle} \hlt{Large Tick Futures}\\
Futures often have economically large ticks. A binding tick creates bid-ask bounce and price reversals, but it also supports large displayed quantities at the best limits relative to comparable equities.
\end{principle}

\begin{taxonomy} \hlt{Futures Architecture Drivers}
\begin{enumerate}[label=\roman*.]
\setlength{\itemsep}{0pt}
\item \hlt{standardisation}: contracts are fungible by expiry and specification;
\item \hlt{central clearing}: margin and default management are built into the market;
\item \hlt{leverage}: small price moves create large economic exposure;
\item \hlt{hedging link}: futures connect to cash baskets, ETFs, options, and swaps.
\end{enumerate}
\end{taxonomy}

\begin{mechanism} \hlt{Variation Margin}\\
Futures gains and losses are settled through daily variation margin.\\
The broker or FCM manages client margin calls, liquidation rights, and cash balances; credit risk is therefore controlled continuously rather than only at final settlement.
\end{mechanism}

\begin{taxonomy} \hlt{Futures Block and EFP Segmentation}
\begin{enumerate}[label=\roman*.]
\setlength{\itemsep}{0pt}
\item \hlt{block trade}: negotiated futures size reported under exchange rules;
\item \hlt{EFP}: exchange cash or physical exposure for offsetting futures;
\item \hlt{transitory EFP}: short-lived EFP used when the book is closed or thin.
\end{enumerate}
They move large risk off-book, subject to eligibility and reporting rules.
\end{taxonomy}
